\section{The inclusion-minimal planar construction}\label{sec:background}

We recall the parts of the planar argument of \citet{BukhJeffs2023} that
are needed for the relative theorem.  The two points that matter are: a
compact convex realization can be shrunk subject to finitely many retained
incidences, and an inclusion-minimal planar realization is polygonal.

\subsection{Constrained inclusion-minimal tuples}

Let $\mathcal X=(X_0,\ldots,X_n)$ be a tuple of compact convex subsets of
$\mathbb R^2$, and let $R\subset\mathbb R^2$ be finite.  We say that a
subtuple $\mathcal Y=(Y_0,\ldots,Y_n)$ is \emph{$R$-admissible} when
\[
 Y_i\subseteq X_i,\qquad Y_i\cap R=X_i\cap R,
 \qquad \code(\mathcal Y)=\code(\mathcal X).
\]

\begin{lemma}[Constrained minimal tuple]\label{lem:minimal}
Assume that $R$ contains a representative of every word of
$\code(\mathcal X)$.  Then every nonempty family of $R$-admissible tuples
contains a coordinatewise inclusion-minimal member.
\end{lemma}

\begin{proof}
Order the admissible tuples by coordinatewise inclusion.  For a descending
chain, take the coordinatewise intersection.  Compactness and convexity are
preserved, as are the incidences on $R$.  Every source word still occurs,
because $R$ will contain a representative for every closed-code word.  If a
word is realized at a point of the chain intersection, then, for each
coordinate absent at that point, choose a member of the chain that already
excludes the point.  One member of the totally ordered chain is contained in
all finitely many chosen members and realizes the same word.  Thus no new
word appears in the intersection.  Zorn's lemma gives a minimal member.
\end{proof}

\subsection{Planar polygonality}

The geometric result of \citet{BukhJeffs2023} is the following constrained
form of their inclusion-minimality proposition.

\begin{lemma}[Planar polygonality]\label{lem:polygonal}
Every member of an $R$-admissible inclusion-minimal tuple in the plane is a
convex polygon, a segment, a point, or the empty set.
\end{lemma}

\begin{proof}[Proof sketch and retained constraints]
The proof is the local planar simplification argument of
\citet{BukhJeffs2023}.  If one member has infinitely many extreme points,
choose an extreme accumulation point outside the finite representative set.
Inside a sufficiently small disk, every incident convex boundary meets the
circle in one connected arc and no unrelated boundary or representative is
encountered.  Replace each incident set locally by the cone from the chosen
point over its circular trace.  Radial projection to the circle shows that
no membership pattern changes, while at least one coordinate shrinks
properly, contradicting inclusion minimality.

The argument tolerates finitely many additional retained points because the
local disk is chosen disjoint from them.  It also tolerates a coordinate
that is a fixed polygon and contains all its vertices in $R$: a local
simplification at a nonvertex boundary point does not change the convex hull
of those vertices.  Segments, points, and the empty set are already
polytopal.  We use no vertex-count estimate in this paper.
\end{proof}

\subsection{Open codes from closed realizations}

Let $\mathcal U=(U_0,\ldots,U_n)$ be bounded open convex sets in the plane
and put $X_i=\overline{U_i}$.  Passing to closures may introduce codewords,
but only on lower-dimensional active intersections.

\begin{lemma}[Closure-only words]\label{lem:closureonly}
If
$c\in\code(X_0,\ldots,X_n)\setminus\code(U_0,\ldots,U_n)$, then
$\bigcap_{i\in c}X_i$ has empty Euclidean interior.
\end{lemma}

\begin{proof}
Assume the active closed intersection has nonempty interior.  Let $q$ realize
$c$ in the closed tuple and choose $p$ in the interior of the active
intersection.  The segment from $p$ to $q$, except possibly at $q$, lies in
that interior.  Since $q$ has positive distance from every inactive closed
set, a nearby point $q'$ on the segment still avoids all inactive sets.  It
lies in $\operatorname{int}X_i=U_i$ for all active $i$, and therefore
realizes $c$ in the open tuple, a contradiction.
\end{proof}

For every nonempty $\eta$ such that
\[
 X_\eta:=\bigcap_{i\in\eta}X_i
\]
is nonempty and has empty interior, choose one affine line $L_\eta$
containing $X_\eta$; convexity in the plane makes this possible.  Denote the
finite set of these lines by $\mathcal L_{\rm deg}$.

\begin{lemma}[Open-code recovery]\label{lem:openrecovery}
Let $\mathcal Y=(Y_0,\ldots,Y_n)$ be a compact convex subtuple of
$\mathcal X=(X_0,\ldots,X_n)$ with the same closed code.  Suppose:
\begin{enumerate}
\item for every $L\in\mathcal L_{\rm deg}$ and every $i$,
      \[
        L\cap\operatorname{int}Y_i=L\cap U_i;
      \]
\item for every word $c\in\code(\mathcal U)$, a retained triangle has
      vertices in $Y_i$ for all $i\in c$ and contains a source point of
      word $c$ in its interior.
\end{enumerate}
Then
\[
 \code(\operatorname{int}Y_0,\ldots,\operatorname{int}Y_n)
 =\code(\mathcal U).
\]
\end{lemma}

\begin{proof}
First prove the inclusion from the new open code to the source open code.
On every line in $\mathcal L_{\rm deg}$, the complete strict membership
pattern is unchanged by hypothesis.  Let
$L=\bigcup_{\eta}L_\eta$, and take $p\notin L$.  Put
$c=\patt_{\mathcal Y}(p)$.  Since $\code(\mathcal Y)=\code(\mathcal X)$,
this word belongs to the closed source code.  By
\cref{lem:closureonly}, any closed-source word absent from the open source
can occur only inside one of the active intersections used to define
$\mathcal L_{\rm deg}$, and hence only on $L$.  Thus $c$ is a source open
word.

The retained triangle for $c$ implies that
$\bigcap_{i\in c}Y_i$ has nonempty interior; choose
$q$ in that interior.  For sufficiently small $\varepsilon>0$, set
\[
 p_\varepsilon=p+\varepsilon(p-q).
\]
The perturbation continues the ray from the interior point $q$ through
$p$.  If $p\in\operatorname{int}Y_i$, then $p_\varepsilon$ remains in
$Y_i$ for all sufficiently small $\varepsilon$.  If
$p\in\partial Y_i$ and $q\in\operatorname{int}Y_i$, convexity implies
that every sufficiently small point beyond $p$ on this ray lies outside
$Y_i$.  Finally, if $p\notin Y_i$, nonmembership persists for small
$\varepsilon$ because $Y_i$ is closed.  Hence
\[
 \patt_{(\operatorname{int}Y_i)}(p)
 =\patt_{\mathcal Y}(p_\varepsilon).
\]
Taking $\varepsilon$ small also keeps $p_\varepsilon\notin L$, so the
right-hand word is a source open word.  This proves that no new open word is
created away from $L$, and the trace hypothesis handles points on $L$.

Conversely, take $c\in\code(\mathcal U)$ and its retained triangle.  The
source witness in the interior of that triangle lies in
$\operatorname{int}Y_i$ for every $i\in c$.  For $j\notin c$, the witness
lies outside $U_j=\operatorname{int}X_j$; since $Y_j\subseteq X_j$, it
cannot lie in $\operatorname{int}Y_j$.  Hence the witness retains the exact
word $c$.
\end{proof}

The preceding proof is the open-to-closed step of
\citet{BukhJeffs2023}, with the finite line traces and witness triangles
made explicit.  The prescribed lines in our theorem will simply be added to
$\mathcal L_{\rm deg}$.
